%
% File: chap01.tex
% Author: Solal Rapaport
% Description: Introduction chapter. Social-engineering thread is dropped from this thesis.
%
\chapter{Introduction}
\label{chap:intro}

\initial{I}n March 2025, attackers altered \texttt{tj-actions/changed-files}, a widely reused GitHub Actions helper, so that \num{346} existing versions --- from \texttt{v39} through \texttt{v47} --- all pointed to one malicious commit~\cite{paloalto_unit42_2025}.
Projects that resolved any of those familiar version names then executed attacker-controlled code capable of exposing developers' credentials.

The incident required no vulnerability in the downstream projects.
Its leverage came from a mismatch between what Git permits and what software consumers expect: anyone with sufficient repository access can retarget a tag, whereas workflows commonly treat a version name as a stable identifier.
The attack was conspicuous once discovered, but the underlying operation was an ordinary reference update.
Within the vast volume of legitimate repository changes, comparable events need not announce themselves as attacks.
Finding them is therefore an automated detection problem akin to searching for a needle in a haystack: absence of a detected incident does not establish that repository state remained unchanged or trustworthy.

This chapter frames that problem at three nested levels.
A visible release name may resolve to different code over time; a published commit history may be replaced; and a file removed during such a replacement may remain available in an archive, possibly with sensitive contents.
The challenge is, besides showing that Git permits these operations, to identify which changes to search for, detect them efficiently across an ecosystem-scale archive, and classify the resulting candidates without treating every maintenance operation as a security incident.

\section{Software Repositories as Load-Bearing Infrastructure}
\label{sec:intro-context}

Open-source repositories are load-bearing infrastructure for web services, embedded systems, scientific computing, and public digital services.
A deployable artifact is commonly assembled from hundreds or thousands of upstream packages, themselves connected by further dependencies.
Its production therefore depends on repository states maintained by people and organisations that downstream developers may never interact with directly.

\label{subsec:intro-oss-critical}%
This dependency structure makes repository integrity a condition of digital trust.
Package managers resolve version specifications, continuous integration systems fetch and execute dependencies, and lockfiles or software bills of materials record inputs for later reproduction and audit.
Frameworks such as SLSA (Supply-chain Levels for Software Artifacts) likewise associate provenance with named source inputs~\cite{tamanna2024analyzingchallengesdeploymentslsa}.
These mechanisms must support both security against malicious substitution and operational safety: a declared input should continue to designate the intended software and guarantee dependable production, reproduction, and audit.

\label{subsec:intro-trust-anchors}%
Cryptographic identifiers can provide a technical form of this guarantee.
A hash is computed from an artifact's bytes: if those bytes change, the hash changes with overwhelming probability.
A consumer that records and verifies the expected hash can therefore detect substitution.
Most human-facing dependency declarations, however, use names such as \texttt{v2.3.1} or \texttt{main} rather than content hashes.
Those names are convenient, but Git does not bind them permanently to one object.

Public repositories consequently act as trust anchors without always providing immutability.
A workflow containing \texttt{tj-actions/changed-files@v45} delegates the meaning of \texttt{v45} to the upstream repository.
A later audit similarly assumes that the recorded version still identifies the source that was compiled~\cite{Vu21}.
If the name moves, the same declaration can yield different software without any local change.

\label{subsec:intro-security-assumptions}%
The motivating question is thus whether changes that undermine this trust can be detected after the fact and at ecosystem scale.
A code forge such as GitHub or GitLab usually presents the current repository state, not the state that a mutable name designated yesterday.
Once a reference or history is replaced, the most direct evidence disappears from that interface.
An empirical investigation needs repeated observations from an independent source, together with methods capable of reducing an enormous stream of ordinary changes to interpretable candidates.

\section{The Gap Between Git's Design and Its Consumption}
\label{sec:intro-problem}

\label{subsec:intro-mutability-challenge}%
Git stores file contents, directory trees, and commits as content-addressed objects~\cite{ProGit2014}.
The identifier of a commit covers its content, metadata, and parent identifiers.
Hence, changing an earlier commit creates a new identifier for that commit and for every descendant that refers to it; this propagation is called the \emph{snowball effect}.
Therefore, an existing object cannot silently change its content and keep the same identifier.

Unlike content-addressed objects, references --branches or tags-- do not have that property.
A reference is a name that points to an object identifier, and Git permits that pointer to be updated or deleted~\cite{ProGit2014,git_tag_retagging_docs}.
The \emph{Pro Git} description of a tag as ``very much like a branch that doesn't change'' is consequently an expectation about use, not a protocol guarantee~\cite{ProGit2014}.
\Cref{chap:background} develops this object--reference distinction in full.

\label{subsec:intro-benign-vs-security}%
This design exposes three different observable signals.
First, the target of a release tag can change while its version name remains in registries, documentation, and workflow files.
Second, a force-push can make previously published commits unreachable from a branch and replace them with a different graph.
Third, files removed by such rewriting may survive in clones and archives, so repository removal alone cannot establish that sensitive contents have been revoked.

\label{subsec:intro-reproducibility-auditing}%
None of these signals should be instantly labeled as incidents.
Repository states reveal what changed, but generally not why it changed.
Maintainers often move tags to repair releases, rewrite histories during routine development, and remove credentials as a genuine remediation step.
Automatic analysis must therefore separate detection from interpretation: first recover observable changes, then classify their form and narrow the cases that warrant security analysis.

\section{Software Heritage as an Empirical Lever}
\label{sec:intro-objectives}

\label{subsec:intro-obj-tags}%
\SWH, the largest public archive of source code, supplies the missing historical observations~\cite{swhcacm2018}.
It visits public code forges and records snapshots of repository state.
It also assigns persistent, content-derived identifiers (SWHIDs) to archived objects, from file contents to complete snapshots~\cite{iso-swhid}.
Because a SWHID is derived from the object itself, two observations with the same identifier denote the same archived artifact, whereas different identifiers make a change explicit.

\label{subsec:intro-obj-histories}%
Two visits to the same origin provide a before-and-after pair.
Comparing consecutive snapshots reveals reference names that appeared or disappeared and names whose targets differ.
The earlier snapshot can retain objects no longer advertised by the forge, allowing later analysis of both sides of an alteration.
This turns otherwise transient repository changes into durable evidence.

\label{subsec:intro-obj-secrets}%
The archive solves the observability problem but creates a computational one.
The tag analysis alone considers \OriginsWithTagsShort repositories over nine years, while the history analysis considers \OriginsTwoVisitsOrMoreShort origins in an archive graph containing \SWHRevisionsShort revisions.
At this scale, exhaustive manual inspection is impossible, while indiscriminate graph traversal, content retrieval, and secrets scanning are unnecessarily expensive.
The empirical problem is to, first, sequence automatic operations so that cheap structural comparisons generate candidates, then use taxonomies to distinguish what changed, and, finally, run progressively more expensive analyses that operate only on relevant subsets.

\section{Three Studies of Repository Integrity}
\label{sec:intro-rqs}

The three studies progressively refine that search.
They move from reference-level differences, through commit-graph differences, to the contents of selected removed files.
Each study examines a deeper layer of repository state.

\subsection{When Tag Names Outlive Their Targets}
\label{subsec:intro-rq-tags}

The first study asks how often a stable tag name is observed with different targets, which forms those changes take, and how they relate to downstream visibility~\cite{rapaport2026tagalterations}.
The challenge is to follow each tag name across irregular archival visits and classify what changed: deletion, recreation, a metadata-only update, or a retarget that points the same name at different source code.
The archive shows what changed, not why, yet that classification still matters: retargeting a tag can be routine release maintenance or the mechanism behind a supply-chain attack, and Git alone does not reveal which.
\Cref{chap:tags} develops the definitions, taxonomy, and evidence needed to make that distinction.

\subsection{Commit Graphs Are Not Append-Only}
\label{subsec:intro-rq-histories}

The second study asks when a branch that appears to advance has instead replaced part of its published history, which repositories and branches are affected, and what kinds of changes caused the replacement~\cite{DBLP:conf/kbse/RapaportPTZ25}.
This is harder than comparing two branch tips.
One direct edit can change the identifiers of a long chain of descendants through the snowball effect, and equivalent branch roles occur under many names.
Automatic detection must traverse compressed commit graphs, isolate root causes from propagated changes, and normalise structural variation before any alteration can be interpreted~\cite{torresarias2016gitmetadata}.
\Cref{chap:histories} addresses that graph-scale detection and classification problem.

\subsection{Removal Is Not Revocation}
\label{subsec:intro-rq-secrets}

The final study asks which files removed through history rewriting are likely to contain credentials and what can responsibly be established about those credentials.
The archive holds too many altered files to retrieve and scan indiscriminately, and no ground truth identifies a credential-bearing file in advance.
Even a broad filename selection produces \SecretFirstExtractionShort records; a more selective pass must reduce them to \SecretSecondExtractionShort candidates before archived contents can be retrieved.
Aggressive filename filtering keeps the study tractable, but it may miss real cases.
Scanners can raise false alarms.
Identical archived files should be scanned once and then traced back to every repository where they appeared.
Finding a secret in the archive does not prove it still grants access.
Chapter~\ref{chap:secrets} develops a staged method to decide which removed files to examine, how to scan them, and what claims automatic analysis can responsibly make.

\section{Cross-Cutting Detection Challenges}
\label{sec:intro-contributions}

\subsection{Candidate Generation at Archive Scale}
\label{subsec:intro-contrib-empirical}

Automatic detection begins with a deliberately broad structural question: did the target of a name change, or did a commit previously reachable from a branch become unreachable?
These tests can be applied to archived identifiers and graph topology without downloading every source file.
Only candidates that survive this first stage justify deeper object comparison, file retrieval, or content scanning.
The order is essential: at archive scale, making an expensive operation slightly more accurate is insufficient if it must still run over every object.

\subsection{Taxonomy and Staged Filtering}
\label{subsec:intro-contrib-methodological}

Candidate generation answers whether an observable difference exists but does not explain the difference, so we introduce taxonomies and filters as additional tools.
Study-specific taxonomies determine whether source content, commit metadata, reference structure, or filenames changed.
For the next stage, repositories need repeated observations, tag analysis needs tag-bearing origins, history analysis need root-cause commits, and secrets analysis needs only relevant filenames.
This is where filters come in, allowing us to conduct our studies on relevant data only.
This staged design makes the computation manageable while keeping explicit which events each filter can exclude.

\section{Thesis Structure}
\label{sec:intro-structure}

\Cref{chap:background} introduces the Git, release-engineering, supply-chain, credential, and \SWH concepts needed to formalise the problem.
\Cref{chap:tags} specifies the snapshot-pair methodology and applies it to release-tag alterations.
\Cref{chap:histories} extends the search from references to commit graphs and develops an operational detector.
\Cref{chap:secrets} narrows file-removal candidates through staged filtering, content scanning, and non-intrusive validation.
\Cref{chap:conclusion} synthesises what these studies imply for repository integrity and for future monitoring.
